\section{Conclusion}

We have constructed a transport system on the flag space of the
dodecahedron, producing a $60$-vertex chamber graph $\Gsixty$ with
successive quotients
\[
\Gsixty \to \Gthirty \to \Gfifteen,
\]
where
\[
\Gfifteen \cong L(\mathrm{Petersen}).
\]

The lifted transport induces a canonical $\ZZtwo$-valued cocycle on
$\Gfifteen$, realized through the double cover
\[
\Gthirty \to \Gfifteen.
\]
This cocycle has minimal support size $6$, cannot be recovered from local
combinatorial data on the core graph, and represents a nontrivial class in
\[
H^1(\Gfifteen,\ZZtwo).
\]
Its natural observable is loop holonomy: nontrivial cycles detect the
failure of lifted transport to close on the same sheet.

From the same transport mechanism we also derive the sector identity
\[
M M^{\mathsf T} = A^3 + 2A^2 + 2I,
\]
relating the sector--edge incidence operator to the adjacency algebra of
$\Gfifteen$. After centering and normalization, this yields a spherical
embedding of $\Gfifteen$ in $\mathbb{R}^{14}$ with exactly three distinct
inner products.

These two structures arise from the same transport system but play
different roles. The cocycle captures a cohomological obstruction visible
through lifted loop holonomy, while the sector system encodes a rigid
quadratic incidence geometry on the core graph.

Taken together, these results show that local transport rules on a
combinatorial polyhedron can produce both nontrivial topological
structure and rigid metric geometry on a canonical quotient graph.

More broadly, this construction exhibits a mechanism by which transport on
flag spaces induces canonical algebraic and geometric structures on core
quotients, linking double covers, holonomy, adjacency algebras, and
spherical configurations in a unified framework.
